% Chapter 1: Introduction — problem motivation, thesis contribution (CascadeLP), roadmap
\selectlanguage{greek}


\chapter{Εισαγωγή}
\label{ch:introduction}


\section{Παρακίνηση}
\label{sec:motivation}

Το σύνολο της ανθρώπινης γνώσης είναι εκ φύσεως οργανωμένο ως γράφημα.
Η
\en{Wikipedia} μόνη της κωδικοποιεί εκατοντάδες χιλιάδες οντότητες---ανθρώπους,
τόπους, οργανισμούς, έννοιες---και οι υπερσύνδεσμοι μεταξύ των αντίστοιχων άρθρων
κωδικοποιούν μια πλούσια σχεσιακή δομή που αντικατοπτρίζει τον τρόπο με τον οποίο
συνδέονται οι ιδέες.
Η πρόβλεψη ποιες ακμές απουσιάζουν από ένα τέτοιο γράφημα, ένα
πρόβλημα γνωστό ως \textit{πρόβλεψη ακμών}, έχει εφαρμογές που εκτείνονται από τη
συμπλήρωση βάσεων γνώσης και τα συστήματα συστάσεων έως την ανάλυση δικτύων
παραπομπών και την ανακάλυψη αλληλεπιδράσεων φαρμάκων.
Η πρόκληση εντείνεται όταν οι
κόμβοι φέρουν πλούσια κειμενικά χαρακτηριστικά: η τοπολογία του γραφήματος από μόνη της
ενδέχεται να είναι αραιή ή ασαφής, ενώ το κείμενο από μόνο του μπορεί να έχει πολύ
υψηλή διαστατικότητα για απλά μέτρα ομοιότητας.
Η από κοινού αξιοποίηση και των δύο
πηγών πληροφορίας είναι επομένως τόσο αναγκαία όσο και μη τετριμμένη.

Οι παραδοσιακές μέθοδοι πρόβλεψης ακμών κατατάσσονται σε δύο ευρείες κατηγορίες.
Οι
δομικοί ευρετικοί δείκτες---όπως οι \en{Common Neighbours}, \en{Adamic-Adar} και
\en{Preferential Attachment}---είναι υπολογιστικά φθηνοί και ερμηνεύσιμοι, αλλά
αγνοούν το κειμενικό περιεχόμενο και αποτυγχάνουν πλήρως όταν δύο κόμβοι δεν
μοιράζονται κανέναν κοινό γείτονα.

Στο αντίθετο
άκρο του φάσματος, τα μεγάλα προ-εκπαιδευμένα γλωσσικά μοντέλα (\en{PLMs}) όπως το
\en{BERT} \cite{devlin2019bert} και οι απόγονοί του μπορούν να κωδικοποιούν πλούσιες
σημασιολογικές σχέσεις απευθείας από κείμενο, και πρόσφατες εργασίες κατέδειξαν ότι
προσαρμοσμένοι κωδικοποιητές προτάσεων μπορούν να επιτύχουν σχεδόν τέλεια ακρίβεια
στο σύνολο αξιολόγησης \en{DSAA 2023} για πρόβλεψη ακμών στη \en{Wikipedia}
\cite{phan2023link,tran2023text}.
Ωστόσο, η κωδικοποίηση κειμένου με μετασχηματιστή αποτελεί το ακριβότερο στάδιο
της σημασιολογικής αλυσίδας, και οι αναφερόμενες σχεδόν τέλειες
βαθμολογίες εγείρουν ερωτήματα ως προς τα χαρακτηριστικά του συνόλου δεδομένων---ιδίως
κατά πόσον ζεύγη που είναι τετριμμένα αναγνωρίσιμα, όπως αυτοβρόχοι ή δομικά
προφανείς σύνδεσμοι, διογκώνουν τη συνολική μετρική.

Η παρούσα εργασία παρακινείται από τρεις παρατηρήσεις που παραμένουν ανεπαρκώς
διερευνημένες στην υπάρχουσα βιβλιογραφία.
Καταρχάς, η εξαγωγή βαθιών σημασιολογικών χαρακτηριστικών είναι σαφώς ακριβότερη
από τον υπολογισμό τοπικών ευρετικών δεικτών. Σε εφαρμογή όπου τα χαρακτηριστικά
παράγονται κατά απαίτηση, η δρομολόγηση μπορεί να περιορίσει το πλήθος των κόμβων
για τους οποίους απαιτείται κωδικοποίηση· όταν οι ενσωματώσεις προϋπολογίζονται και
επαναχρησιμοποιούνται, το μετρούμενο όφελος αφορά το στάδιο απόφασης και όχι την
ήδη καταβληθείσα κωδικοποίηση.
Επιπλέον, διαφορετικοί υπο-πληθυσμοί ζευγών κόμβων
ανταποκρίνονται πολύ διαφορετικά σε διαφορετικές οικογένειες χαρακτηριστικών: τα
δομικά πλούσια ζεύγη ταξινομούνται τετριμμένα από την τοπολογία, τα ζεύγη \en{cold-start}
με μηδενικούς κοινούς γείτονες απαιτούν σημασιολογικά χαρακτηριστικά,
και τα πραγματικά αμφίβολα ζεύγη ενδέχεται να χρειαστούν βαθιά γλωσσική κατανόηση.
Τέλος, οι σχεδόν τέλειες βαθμολογίες που αναφέρονται από τους \textcite{phan2023link}
και \textcite{tran2023text} εγείρουν ένα ερώτημα που η υπάρχουσα
βιβλιογραφία δεν απαντά ρητά: οφείλεται η απόδοση αυτή σε πραγματική σημασιολογική
συλλογιστική, ή στο ότι το σύνολο αξιολόγησης περιέχει ένα μεγάλο κλάσμα ζευγών που
είναι τετριμμένα ταξινομήσιμα (αυτοβρόχοι, ζεύγη με πολλούς κοινούς γείτονες) πριν
καν κληθεί οποιοδήποτε μοντέλο; Ένας αρχικός έλεγχος επικάλυψης που πραγματοποιήσαμε
στο \texttt{\en{train.csv}}/\texttt{\en{test.csv}} του \en{DSAA 2023} βρίσκει
αμελητέα επικάλυψη ζευγών (\LeakageExact{} ακριβή ζεύγη και \LeakageReversed{} αντεστραμμένα ζεύγη σε \TestPairs{}
ζεύγη ελέγχου). Ο έλεγχος αυτός αποκλείει την απλή απομνημόνευση
επαναλαμβανόμενων ζευγών, όχι όμως άλλες αστοχίες του πρωτοκόλλου ή της
δειγματοληψίας. Πράγματι, η μεταγενέστερη ανάλυση αποκαλύπτει ότι ένας πολύ μικρός
αριθμός τιμών του πρώτου άκρου καθορίζει σχεδόν πλήρως την ετικέτα, καθιστώντας τη
βαθμολογία του \en{DSAA 2023} ανεπαρκή ως απόδειξη γνήσιας σχεσιακής πρόβλεψης.

Χρησιμοποιούμε το \en{DSAA 2023} ως πρώτο πεδίο αξιολόγησης επειδή ικανοποιεί τρεις
προϋποθέσεις που σπάνια συνυπάρχουν σε ένα μόνο σύνολο δεδομένων: πρόκειται για
πραγματικό γράφημα με κειμενικά χαρακτηριστικά κόμβων σε κλίμακα εκατοντάδων
χιλιάδων κόμβων και όχι για συνθετικό ή τεχνητά μικρό στιγμιότυπο, διαθέτει
δημόσια ετικέτες εκπαίδευσης και ελέγχου που επιτρέπουν αναπαραγώγιμη μέτρηση, και
συνοδεύεται από έναν καθιερωμένο πίνακα κατάταξης διαγωνισμού με δεκάδες υπαρκτές
λύσεις προς σύγκριση, αντί για μια μεμονωμένη δημοσιευμένη βαθμολογία χωρίς πλαίσιο
αναφοράς. Επειδή όμως ο έλεγχος του συνόλου αποκάλυψε σοβαρό τεχνούργημα
δειγματοληψίας αρνητικών, κατασκευάζουμε και δεύτερο σύνολο από πραγματικούς
υπερσυνδέσμους της Απλής Αγγλικής \en{Wikipedia}. Η δεύτερη αξιολόγηση δεν αποτελεί
μεταφορά μάθησης: το ίδιο αρχιτεκτονικό σχήμα επανεκπαιδεύεται από την αρχή και
ελέγχεται με ομαδοποιημένη διαίρεση και απόκρυψη της ακμής-στόχου.

Οι παρατηρήσεις αυτές παρακινούν μια \textit{πολυεπίπεδη} προσέγγιση: εφαρμογή
πρώτα του υπολογιστικά φθηνότερου ταξινομητή, κλιμάκωση σε ακριβότερους ταξινομητές
μόνο όταν το φθηνότερο επίπεδο εκδηλώνει αβεβαιότητα, και μέτρηση τόσο της ακρίβειας
όσο και του κόστους συμπερασμού σε κάθε επίπεδο.
Μια τέτοια αλυσίδα κλιμάκωσης
ευθυγραμμίζεται καλά με το πρακτικό σενάριο ανάπτυξης όπου ένα μεγάλο κλάσμα ζευγών
μπορεί να επιλυθεί φθηνά, και μόνο ένα μικρό υπόλοιπο απαιτεί βαριά υπολογιστική
επεξεργασία.


\section{Ορισμός Προβλήματος}
\label{sec:problem-definition}

Έστω $G = (V, E)$ ένα μη κατευθυνόμενο γράφημα όπου κάθε κόμβος $v \in V$ συνδέεται
με ένα κειμενικό έγγραφο $d_v$ (στο πλαίσιο της παρούσας εργασίας, το κείμενο του
άρθρου της \en{Wikipedia}). Δεδομένου ενός συνόλου ζευγών κόμβων $\mathcal{Q} =
\{(u, v)\}$, το πρόβλημα της \textit{πρόβλεψης ακμών} συνίσταται στην απόδοση μιας
δυαδικής ετικέτας $y_{uv} \in \{0, 1\}$ σε κάθε ζεύγος, όπου $y_{uv} = 1$ δηλώνει
ότι μια ακμή $(u, v) \in E$ υπάρχει (ή πρέπει να υπάρχει) και $y_{uv} = 0$ δηλώνει
απουσία ακμής.
Το πρόβλημα αξιολογείται σε υπολειπόμενα ζεύγη των οποίων οι
πραγματικές ετικέτες παρακρατούνται κατά τον χρόνο συμπερασμού.

Το πρόβλημα αυτό είναι δύσκολο σε γενικό επίπεδο για δύο συμπληρωματικούς λόγους.
Πρώτον, η τοπολογία του γραφήματος από μόνη της δεν είναι πάντα πληροφοριακή: κόμβοι
που μόλις προστέθηκαν, ή που ανήκουν σε αραιά συνδεδεμένες συνιστώσες, δεν έχουν
κοινούς γείτονες με κανέναν υποψήφιο σύντροφο, οπότε κάθε ευρετικός δείκτης που
βασίζεται σε γειτονιά (\en{Common Neighbours}, \en{Adamic-Adar}) επιστρέφει μηδέν
ανεξαρτήτως της πραγματικής ετικέτας.
Δεύτερον, η κειμενική ομοιότητα δεν είναι
αξιόπιστη υποκατάστατη μεταβλητή για την ύπαρξη ακμής: δύο κόμβοι μπορεί να είναι
θεματικά όμοιοι χωρίς να συνδέονται (π.χ. δύο άρθρα για διαφορετικά είδη με σχεδόν
πανομοιότυπη δομή \en{infobox}) ή να συνδέονται χωρίς να μοιράζονται σχεδόν καθόλου
λεξιλόγιο.
Ένα γενικό σύστημα πρέπει επομένως να αποφασίζει, ανά ζεύγος, ποια πηγή
πληροφορίας---τοπολογία, κείμενο, ή και οι δύο---είναι αξιόπιστη, χωρίς να γνωρίζει
εκ των προτέρων σε ποιο καθεστώς ανήκει το ζεύγος.
Μια \textit{γενική λύση}, με την
έννοια που χρησιμοποιείται στην παρούσα εργασία, είναι μια αρχιτεκτονική που
παραμένει ανεξάρτητη από τις ιδιαιτερότητες ενός συγκεκριμένου συνόλου δεδομένων:
δεν προϋποθέτει συγκεκριμένη γλώσσα, δομή \en{infobox}, ή κατανομή θετικών και
αρνητικών ζευγών, και τα όριά της ρυθμίζονται εμπειρικά σε οποιοδήποτε νέο σύνολο
δεδομένων που ακολουθεί το ίδιο γενικό σχήμα: γράφημα, κείμενο ανά κόμβο, ετικέτες
ζευγών.

Το γενικό πρόβλημα εμφανίζεται σε πολλές παραλλαγές, ανάλογα με το ποια πηγή
πληροφορίας είναι διαθέσιμη ή αξιόπιστη για ένα δεδομένο ζεύγος.
Στην περίπτωση
\en{cold-start}, ένας ή και οι δύο κόμβοι δεν έχουν ακμές στο γράφημα εκπαίδευσης---
είτε επειδή προστέθηκαν πρόσφατα είτε επειδή ανήκουν σε αραιά συνδεδεμένη
συνιστώσα---οπότε κάθε δομικός ευρετικός δείκτης είναι εξ ορισμού αδιάφορος και η
απόφαση πρέπει να βασιστεί αποκλειστικά στο κείμενο.
Στην αντίθετη περίπτωση, όπου
η τοπολογία είναι πυκνή αλλά το κείμενο ελλιπές ή θορυβώδες, ισχύει το αντίστροφο.
Ανάμεσα στα δύο άκρα βρίσκονται τα ζεύγη όπου και οι δύο πηγές πληροφορίας
διατίθενται αλλά διαφωνούν, οπότε το σύστημα πρέπει να αποφασίσει ποια να
εμπιστευτεί.
Κάθε μία από αυτές τις παραλλαγές συνεπάγεται διαφορετικό
συμβιβασμό κόστους-ακρίβειας: οι φθηνοί ευρετικοί δείκτες αρκούν όταν η
τοπολογία είναι πληροφοριακή, ενώ η κλιμάκωση σε ακριβότερα σημασιολογικά
μοντέλα είναι αναγκαία μόνο όταν δεν είναι---και ένα γενικό πλαίσιο πρέπει να
αναγνωρίζει ανά ζεύγος σε ποιο καθεστώς βρίσκεται, χωρίς να το γνωρίζει εκ των
προτέρων.

Στο σύνολο αξιολόγησης \en{Wikipedia DSAA 2023} \cite{phan2023link}, το γράφημα $G$
έχει \GraphNodesTotal{} κόμβους που αντιστοιχούν σε άρθρα της \en{Wikipedia}, εκ των οποίων
\GraphNodes{} εμφανίζονται στη δομική (γειτνιακή) συνιστώσα.
Ετικέτες εκπαίδευσης
παρέχονται για \TrainPairs{} ζεύγη κόμβων (\TrainPosPct\% θετικά, \TrainNegPct\% αρνητικά), και η μετρική
αξιολόγησης είναι η βαθμολογία \en{Macro F1} που υπολογίζεται επί των δυαδικών
θετικών και αρνητικών κλάσεων.
Το σύνολο ελέγχου περιέχει \TestPairs{} ζεύγη.
Το κείμενο
των κόμβων είναι διαθέσιμο ως αρχείο \texttt{\en{nodes.tsv}} που περιέχει τίτλους,
περιγραφές και προαιρετικά περιεχόμενο \en{infobox} της \en{Wikipedia}.

Το γενικό πρόβλημα που αντιμετωπίζει η παρούσα εργασία είναι η κατασκευή ενός
πλαισίου πρόβλεψης ακμών για γραφήματα με κειμενικά χαρακτηριστικά κόμβων που (α)
αξιοποιεί από κοινού τη δομή του γραφήματος και το κείμενο των κόμβων για την
παραγωγή ακριβών προβλέψεων, (β) διαχειρίζεται ρητά το κόστος συμπερασμού
δρομολογώντας εύκολα ζεύγη μακριά από ακριβά μοντέλα, (γ) είναι διαφανές ως προς
το ποια οικογένεια χαρακτηριστικών είναι υπεύθυνη για κάθε πρόβλεψη, και (δ)
χειρίζεται ζεύγη \en{cold-start} όπου ένας ή και οι δύο κόμβοι δεν έχουν
δομικούς γείτονες στο $G$.
Το σύνολο αξιολόγησης \en{Wikipedia DSAA 2023} χρησιμοποιείται στην παρούσα
εργασία ως μία συγκεκριμένη υλοποίηση αυτού του γενικού προβλήματος, όχι ως ο
τελικός στόχος του: παρέχει ένα πραγματικό γράφημα κλίμακας εκατοντάδων χιλιάδων
κόμβων, κείμενο ανά κόμβο και ετικέτες εκπαίδευσης/ελέγχου που καθιστούν το
πλαίσιο μετρήσιμο, χωρίς όμως καμία από τις σχεδιαστικές αποφάσεις του πλαισίου
να προϋποθέτει τις ιδιαιτερότητές του.
Το πλαίσιο πρέπει να λειτουργεί κατά τον χρόνο
συμπερασμού χωρίς πρόσβαση στις πραγματικές ετικέτες των ζευγών ελέγχου,
χρησιμοποιώντας αποκλειστικά το γράφημα εκπαίδευσης, το κείμενο των κόμβων και τις
ετικέτες ζευγών εκπαίδευσης.

Παράλληλος στόχος είναι ο έλεγχος της εγκυρότητας των δύο πεδίων αξιολόγησης:
εντοπίζουμε επαναλαμβανόμενα ζεύγη, αυτοβρόχους, πηγές με σχεδόν σταθερή ετικέτα
και χαρακτηριστικά που διαχωρίζουν εύκολα τις κλάσεις. Η ανάλυση αυτή δεν
ταυτίζει όσα ζεύγη αποτυγχάνουν στα απλά κριτήρια με ζεύγη που απαιτούν κατ'
ανάγκην σημασιολογική συλλογιστική· τα χρησιμοποιεί ως λειτουργικές κατηγορίες
δυσκολίας για τη σύγκριση της συμπεριφοράς των επιπέδων.


\section{Συνεισφορές}
\label{sec:contributions}

Οι κύριες συνεισφορές της παρούσας διπλωματικής εργασίας είναι οι εξής.

\begin{enumerate}

    \item \textbf{\en{CascadeLP}: μια αρχιτεκτονική κλιμάκωσης βάσει αξιοπιστίας για
    πρόβλεψη ακμών με επίγνωση κόστους.}
    Σχεδιάζουμε και υλοποιούμε μια τετραεπίπεδη αλυσίδα κλιμάκωσης
    (αυτοβρόχοι, δομικοί ευρετικοί δείκτες, χαρακτηριστικά μερών του λόγου,
    ομοιότητα \en{Sentence-Transformer}) που δρομολογεί κάθε ζεύγος κόμβων στο
    φθηνότερο μοντέλο ικανό να το επιλύσει με αξιοπιστία, κλιμακώνοντας στο
    ακριβότερο επόμενο επίπεδο μόνο όταν χρειάζεται.
    Στο ανεξάρτητα κατασκευασμένο σύνολο δεδομένων, το \en{CascadeLP} πετυχαίνει
    \en{Macro F1} = \WfCascadeFone{}, ξεπερνά κάθε μεμονωμένη οικογένεια
    χαρακτηριστικών και επιλύει το \WfTierOnePct\% των ζευγών στο φθηνό δομικό
    επίπεδο. Στο \en{DSAA 2023} πετυχαίνει \en{Macro F1} = \CascadeValFone{} στο
    σύνολο επικύρωσης και \KaggleBaselinePublic{} στο δημόσιο \en{leaderboard},
    αλλά οι τιμές αυτές ερμηνεύονται υπό το πρίσμα του τεχνουργήματος
    δειγματοληψίας που τεκμηριώνεται στην εργασία.
    Η ίδια η αρχιτεκτονική δεν προϋποθέτει τίποτα ειδικό για τη \en{Wikipedia}---
    κάθε επίπεδο δέχεται ως είσοδο ένα ζεύγος κόμβων με προαιρετικό κείμενο και
    προαιρετική θέση σε γράφημα---επομένως γενικεύεται σε κάθε πρόβλημα
    πρόβλεψης ακμών όπου το κόστος συμπερασμού διαφέρει σημαντικά μεταξύ
    οικογενειών χαρακτηριστικών, όχι μόνο σε αυτό το συγκεκριμένο σύνολο
    δεδομένων.
    Για κάθε ζεύγος καταγράφουμε επιπλέον ποιο επίπεδο το επέλυσε,
    ώστε η επίδοση να μπορεί να αναλυθεί ανά επίπεδο και όχι μόνο συνολικά.

    \item \textbf{Έλεγχος διαρροής δεδομένων στο σύνολο αξιολόγησης \en{DSAA 2023}.}
    Πριν εμπιστευτούμε τη βαθμολογία του \en{CascadeLP}, εξετάζουμε το
    \texttt{\en{train.csv}} και το \texttt{\en{test.csv}} για
    ακριβή και αντεστραμμένη επικάλυψη ζευγών $(u,v)$, καθώς και για διπλότυπα
    ζεύγη εντός του \texttt{\en{train.csv}} που θα μπορούσαν να διαρρεύσουν
    μέσω της τυχαίας διαίρεσης εκπαίδευσης-επαλήθευσης.
    Βρίσκουμε αμελητέα
    επικάλυψη (\LeakageExact{} ακριβή και \LeakageReversed{} αντεστραμμένα ζεύγη σε \TestPairs{} ζεύγη ελέγχου),
    αποκλείοντας την απλή επικάλυψη ζευγών ως εξήγηση της υψηλής βαθμολογίας.
    Το πρωτόκολλο ελέγχου δεν εξαρτάται από ιδιαιτερότητες του
    \en{DSAA 2023}---απαιτεί μόνο δύο αρχεία ζευγών με κοινό χώρο κόμβων---και
    μπορεί συνεπώς να εφαρμοστεί αυτούσιο σε οποιοδήποτε σύνολο δεδομένων
    πρόβλεψης ακμών όπου τα σύνολα εκπαίδευσης και ελέγχου προέρχονται από τον
    ίδιο γράφημα, κάτι που το καθιστά έναν γενικό έλεγχο ακεραιότητας και όχι μια
    εφάπαξ επαλήθευση ειδική για το παρόν σύνολο δεδομένων.

    \item \textbf{Έλεγχος του \en{DSAA 2023} και εντοπισμός τεχνουργήματος
    δειγματοληψίας αρνητικών.} Εντοπίζουμε
    \SelfLoopsTrain{} ζεύγη αυτοβρόχων στο σύνολο εκπαίδευσης (εκ των οποίων \SelfLoopsTrainPos{} θετικά
    και \SelfLoopsTrainNeg{} αρνητικό --- μια ανωμαλία ετικέτας που εξετάζεται στο
    Κεφάλαιο~\ref{ch:analysis}) και \SelfLoopsTest{} στο σύνολο ελέγχου, αναλύουμε την
    κατανομή πλήθους κοινών γειτόνων και ομοιότητας κειμένου σε θετικά και
    αρνητικά ζεύγη εκπαίδευσης, και ορίζουμε ένα κριτήριο «τετριμμένου ζεύγους»
    για την ποσοτικοποίηση του κλάσματος του συνόλου δεδομένων που είναι
    τετριμμένα ταξινομήσιμο προτού κληθεί οποιοδήποτε μοντέλο.
    Επιπλέον, δείχνουμε ότι \HubCount{} τιμές του \texttt{\en{id1}} καλύπτουν το
    \HubRowPct\% των γραμμών και είναι σχεδόν πάντοτε αρνητικές. Ένας απλός
    κανόνας αναζήτησης βάσει αυτών των τιμών πετυχαίνει \en{Macro F1} =
    \HubLookupFone{}, ενώ συμφωνεί με τις προβλέψεις του \en{CascadeLP} στο
    \HubValAgreementPct\% του συνόλου επικύρωσης. Συνεπώς, οι σχεδόν τέλειες
    επιδόσεις που έχουν αναφερθεί στο \en{DSAA 2023} παραμένουν έγκυρες ως
    βαθμολογίες του διαγωνισμού, αλλά δεν τεκμηριώνουν από μόνες τους ικανότητα
    σχεσιακής πρόβλεψης.

    \item \textbf{Μοντέλο αναφοράς κλασικής μηχανικής μάθησης (\en{TF-IDF} + \en{SVM})~\cite{cortes1995support}.}
    Συμπεριλαμβάνουμε ένα \en{SVM} πυρήνα ακτινικής βάσης εκπαιδευμένο σε
    ομοιότητα συνημίτονου \en{TF-IDF} πάνω σε στρωματοποιημένο δείγμα,
    ακολουθώντας το ιστορικό προηγούμενο των \textcite{alhasan2006link}, ως σημείο αναφοράς για το πόσο απαραίτητα είναι τα
    προ-εκπαιδευμένα γλωσσικά μοντέλα σε αυτό το πρόβλημα.
    Το σημείο αυτό δεν αποσκοπεί απλώς σε σύγκριση βαθμολογιών με το
    \en{CascadeLP}, αλλά χρησιμεύει ως ένδειξη του πόσο νωρίς στην κλίμακα
    πολυπλοκότητας μοντέλου κορεννύεται η απόδοση σε προβλήματα πρόβλεψης
    ακμών με κειμενικά χαρακτηριστικά κόμβων---ένα ερώτημα με νόημα ανεξαρτήτως
    του συγκεκριμένου συνόλου δεδομένων που το θέτει.

    \item \textbf{Αξιολόγηση σε δεύτερο σύνολο πραγματικών υπερσυνδέσμων και
    αναπαραγώγιμη υλοποίηση.} Κατασκευάζουμε ανεξάρτητο σύνολο δεδομένων από
    επίσημα αρχεία \en{SQL} της Απλής Αγγλικής \en{Wikipedia}, με πραγματικές
    ακμές και τυχαία μη συνδεδεμένα ζεύγη, και επανεκπαιδεύουμε το ίδιο
    \en{CascadeLP} χωρίς αλλαγή αρχιτεκτονικής. Ολόκληρη η αλυσίδα επεξεργασίας---έλεγχος
    διαρροής, χαρακτηρισμός διαχωρισιμότητας, εξαγωγή χαρακτηριστικών,
    εκπαίδευση μοντέλων, δρομολόγηση αλυσίδας κλιμάκωσης και
    αξιολόγηση---υλοποιείται σε \en{Python} και καθίσταται δημόσια διαθέσιμη.
    Επειδή κάθε στάδιο της αλυσίδας διαβάζει τα δεδομένα εισόδου του από ένα
    ενιαίο σχήμα (γράφημα, κείμενο ανά κόμβο, ετικέτες ζευγών) και όχι από
    κώδικα ειδικό για το \en{DSAA 2023}, η ίδια υλοποίηση μπορεί να
    επαναχρησιμοποιηθεί σε ένα διαφορετικό σύνολο δεδομένων με ελάχιστες
    προσαρμογές---προϋπόθεση απαραίτητη ώστε η μεθοδολογία να θεωρείται γενική
    και όχι απλώς μια λύση προσαρμοσμένη σε αυτό το σύνολο δεδομένων.

\end{enumerate}

Η αξιολόγηση στο δεύτερο σύνολο ελέγχει εάν η αρχιτεκτονική λειτουργεί σε
διαφορετική κατασκευή δεδομένων. Δεν ελέγχει μεταφορά χωρίς επανεκπαίδευση ούτε
γενίκευση εκτός του οικοσυστήματος της \en{Wikipedia}, τα οποία παραμένουν
κατευθύνσεις μελλοντικής εργασίας.


\section{Δομή της Εργασίας}
\label{sec:structure}

Το υπόλοιπο της εργασίας οργανώνεται ως εξής.

Το Κεφάλαιο~\ref{ch:background} επισκοπεί τη σχετική βιβλιογραφία στις τρεις κύριες
κατηγορίες μεθόδων πρόβλεψης ακμών σε γραφήματα με κειμενικά χαρακτηριστικά: δομικοί
ευρετικοί δείκτες και ενσωματώσεις γραφημάτων, προ-εκπαιδευμένα γλωσσικά μοντέλα και
κωδικοποιητές προτάσεων, και υβριδικές αρχιτεκτονικές που συνδυάζουν τοπολογία γραφήματος
με κείμενο.
Επίσης περιγράφει λεπτομερώς το σύνολο αξιολόγησης \en{Wikipedia DSAA
2023} και εντοπίζει τα κενά στην υπάρχουσα εργασία που παρακινούν το \en{CascadeLP}.
Το κεφάλαιο στηρίζεται σε θεμελιώδεις εργασίες που περιλαμβάνουν τον δείκτη
\en{Adamic-Adar} \cite{adamic2003friends}, το μοντέλο προτιμησιακής σύνδεσης
\en{Barabási-Albert} \cite{barabasi1999emergence}, το \en{BERT} \cite{devlin2019bert}, το \en{RoBERTa}
\cite{liu2019roberta}, το \en{Sentence-BERT} \cite{reimers2019sentence}, το
\en{LinkBERT} \cite{yasunaga2022linkbert}, τα \en{GraphFormers}
\cite{yang2021graphformers} και το \en{Patton} \cite{jin2023patton}.

Το Κεφάλαιο~\ref{ch:methodology} παρουσιάζει το πλαίσιο \en{CascadeLP} στο πλήρες
τεχνικό του βάθος.
Καλύπτει την αλυσίδα προεπεξεργασίας δεδομένων (φόρτωση και καθαρισμός
κειμένου), τη μηχανική χαρακτηριστικών για
κάθε επίπεδο (ευρετικοί δείκτες, διανύσματα συχνότητας μερών του λόγου,
ομοιότητα συνημίτονου \en{Sentence-Transformer}), την εκπαίδευση του ταξινομητή κάθε
επιπέδου, τον αλγόριθμο δρομολόγησης βάσει αξιοπιστίας και τη στρατηγική χειρισμού
του \en{cold-start}.

Το Κεφάλαιο~\ref{ch:experiments} περιγράφει τη διάταξη πειραμάτων: τη στρατηγική διαίρεσης
εκπαίδευσης-επαλήθευσης, τη ρύθμιση υπερπαραμέτρων, το πρωτόκολλο αξιολόγησης και
τα μοντέλα αναφοράς που χρησιμοποιούνται για σύγκριση.
Παρουσιάζει τα δύο σύνολα δεδομένων, τα αντίστοιχα πρωτόκολλα και τα κύρια
ποσοτικά αποτελέσματα του \en{CascadeLP} και των μοντέλων αναφοράς.

Το Κεφάλαιο~\ref{ch:analysis} ερμηνεύει συγκριτικά τα αποτελέσματα, τεκμηριώνει
το τεχνούργημα δειγματοληψίας του \en{DSAA 2023}, παρουσιάζει μελέτες
\en{ablation}, ανάλυση των υποσυνόλων χωρίς κοινό γείτονα και των σφαλμάτων, και
διαχωρίζει το μετρούμενο κόστος δρομολόγησης από τη δυνητική εξοικονόμηση
εξαγωγής χαρακτηριστικών.

Το Κεφάλαιο~\ref{ch:conclusion} συνοψίζει τα ευρήματα, συζητά τους περιορισμούς του τρέχοντος
πλαισίου και περιγράφει κατευθύνσεις για μελλοντική έρευνα, συμπεριλαμβανομένης
δυναμικής προσαρμογής ορίων, ενσωμάτωσης νευρωνικών δικτύων γραφημάτων και αξιολόγησης
σε πρόσθετα σύνολα αξιολόγησης γραφημάτων με κειμενικά χαρακτηριστικά.
